\section{Background}
\label{sec:background}

This section summarizes the neuroscience and Buddhist psychology foundations
that inform \alaya{}'s design. A comprehensive treatment of these traditions
and their convergence appears in the companion survey~\cite{hui2026taxonomy};
here we focus on the specific principles that map to architectural decisions.

\subsection{Complementary Learning Systems}

McClelland~\etal~\cite{mcclelland1995} proposed that biological memory
relies on two complementary systems: a fast-learning hippocampal system that
encodes episodes with high fidelity, and a slow-learning neocortical system
that gradually extracts stable patterns through interleaved replay. This
separation prevents catastrophic interference---new learning destroying old
knowledge---while enabling generalization from accumulated experience.

\alaya{} instantiates this with an episodic store (rapid, high-fidelity
conversation recording) and a semantic store (stable, deduplicated knowledge
extracted during consolidation). The consolidation operation implements
CLS-style interleaved replay: batches of recent episodes are processed by an
LLM provider to extract facts, relationships, events, and concepts, which
are stored as semantic nodes distinct in kind from their source episodes.

\subsection{Bjork Dual-Strength Forgetting}

Bjork and Bjork~\cite{bjork1992} distinguished two independent strengths for
every memory: \emph{storage strength} (how deeply encoded) and
\emph{retrieval strength} (how currently accessible). Storage strength
increases monotonically with exposure; retrieval strength decays over time
but recovers dramatically when a memory is successfully retrieved---the
\emph{spacing effect}. Anderson~\etal~\cite{anderson1994} extended this with
\emph{retrieval-induced forgetting} (RIF): retrieving one memory actively
suppresses competitors sharing the same retrieval cues.

\alaya{} implements both mechanisms. Each memory node carries independent
storage and retrieval strengths. Retrieval boosts both strengths; the passage
of time decays only retrieval strength. Competitors sharing retrieval cues
are suppressed during query, implementing RIF. The \texttt{forget()} operation
uses retrieval strength to determine which memories are effectively
inaccessible.

\subsection{Hebbian Learning and Spreading Activation}

Hebb's principle~\cite{hebb1949}---``neurons that fire together wire
together''---describes use-dependent synaptic strengthening. Collins and
Loftus~\cite{collins1975} formalized spreading activation in semantic
networks: activating one concept propagates to associated concepts through
weighted links with distance-based decay.

\alaya{} maintains a Hebbian graph overlay across all memory stores. Links
are created during encoding (temporal, topical, entity, causal) and during
retrieval (co-retrieval links between jointly accessed memories). Link
weights strengthen through repeated co-activation and decay through disuse,
producing emergent small-world topology. Spreading activation is used during
retrieval to discover indirect associations---a query about ``cooking'' can
surface chemistry-related memories through shared Hebbian links.

\subsection{Yogac\=ara Process Model}

The Yogac\=ara school's theory of \=alaya-vij\~n\=ana (storehouse
consciousness) describes memory as a dynamic substrate holding \emph{seeds}
(\emph{b\=ija}) that:
\begin{itemize}[nosep]
    \item \textbf{Strengthen through perfuming} (\emph{v\=asan\=a}): each
    experience leaves a subtle imprint; preferences emerge from accumulated
    impressions, not single observations.
    \item \textbf{Ripen heterogeneously} (\emph{vip\=aka}): seeds mature
    into forms qualitatively different from their causes, paralleling CLS
    consolidation.
    \item \textbf{Transform structurally}
    (\emph{\=a\'sraya-par\=av\d{r}tti}): the store itself reorganizes
    periodically, resolving contradictions and pruning obsolete content.
    \item \textbf{Condition mutually}: seeds activate related seeds,
    paralleling spreading activation in Hebbian networks.
\end{itemize}

\noindent These four principles map directly to \alaya{}'s lifecycle
operations, as detailed in \Cref{sec:architecture}.
